%% CV - Dalfin Rafael Solihin — BCG Working Student GenAI & Workflow Automation
%% Compile with: lualatex main_bcg_genai.tex

\documentclass[11pt,a4paper,sans]{moderncv}
\moderncvstyle{banking}
\moderncvcolor{blue}

% Force the name and section headings to render in moderncv blue (color1).
\renewcommand*{\namefont}{\fontsize{34}{36}\bfseries\upshape}
\colorlet{firstnamecolor}{color1}
\colorlet{lastnamecolor}{color1}
\colorlet{namecolor}{color1}
\renewcommand*{\sectionstyle}[1]{{\sectionfont\color{color1}#1}}

\usepackage[utf8]{inputenc}
\AtEndPreamble{\hypersetup{
    colorlinks=true,
    linkcolor=blue,
    filecolor=magenta,
    urlcolor=blue,
    pdftitle={Dalfin Rafael Solihin - CV},
    pdfpagemode=UseNone,
}}
\usepackage[scale=0.77]{geometry}
\usepackage{import}
\usepackage{needspace}

\name{Dalfin Rafael}{Solihin}
\address{Leipartstr. 28, 81369 Munich, Germany}{}{}
\phone[mobile]{+49 155 664 752 60}
\email{dalfinrafa@gmail.com}
\extrainfo{\href{https://www.linkedin.com/in/dalfin-rafael-solihin}{LinkedIn}, \href{https://github.com/discretediode}{GitHub}}

\begin{document}

\makecvtitle

% ============================================================
%     PROFILE STATEMENT
% ============================================================

\vspace{6pt}
\small{Mechatronics Engineering student at Munich University of Applied Sciences (HM Munich) with hands-on experience building end-to-end LLM automation workflows. Developed an AI CV Automator integrating the Google Gemini API with a secure data pipeline, prompt engineering, and automated PDF assembly --- using Claude Code as the daily agentic development environment. Complementary experience in ML model evaluation (YOLOv8) and Python-based data analysis in a university research setting. Seeking a working student role at BCG to help consulting teams adopt, build on, and extend generative AI solutions.}

% ============================================================
%     CORE COMPETENCIES
% ============================================================

\section{Core Competencies}
\vspace{1pt}
\begin{itemize}

\item \textbf{Generative AI \& LLM Integration}: Google Gemini API, Claude Code (Anthropic), prompt engineering, LLM output evaluation, agentic workflow design

\item \textbf{Workflow Automation \& API Engineering}: End-to-end pipeline design, PII-safe data handling, automated PDF generation (xhtml2pdf), REST API integration

\item \textbf{Python Development}: Data processing scripts, automation tooling, ML experimentation (YOLOv8/Ultralytics), structured technical reporting

\item \textbf{Analytical Problem Solving}: ML metrics analysis, SFRA measurement data analysis, systematic debugging, iterative experimentation in research environments

\item \textbf{Cross-functional Technical Communication}: Translating technical findings for mixed-expertise audiences; comfortable working in German and English with both technical and business stakeholders

\end{itemize}

% ============================================================
%     PROJECT EXPERIENCE
% ============================================================

\section{Project Experience}
\vspace{3pt}
\begin{itemize}

\needspace{5\baselineskip}
\item{\cventry{2025 - present}{AI Application Generator (CV Automator)}{Personal Project}{Python, Google Gemini API, LLM Pipelines}{}{\vspace{1pt}
\begin{itemize}
    \item Built a local LLM automation tool using the Google Gemini API and Python to generate tailored application documents from an Obsidian knowledge vault, developed end-to-end with Claude Code as the agentic coding environment
    \item Engineered a secure data pipeline with PII stripping and re-injection to prevent sensitive personal data from entering the LLM context, then automated PDF assembly via xhtml2pdf
    \item Designed and iterated prompt templates; evaluated LLM output quality across model versions, demonstrating practical prompt engineering and generative AI product thinking directly applicable to BCG's AI enablement work
\end{itemize}}}

\vspace{3pt}

\needspace{5\baselineskip}
\item{\cventry{2024}{YOLOv8 Object Detection Model}{Personal Project}{Python, YOLOv8/Ultralytics, Computer Vision}{}{\vspace{1pt}
\begin{itemize}
    \item Trained and evaluated a custom YOLOv8 model for image-based object classification, covering dataset preparation, training configuration, and quantitative performance evaluation
    \item Developed familiarity with the full ML experimentation cycle: data preparation, model training, metrics analysis, and iterative improvement --- skills directly applicable to evaluating new AI tools and features
\end{itemize}}}

\vspace{3pt}

\needspace{5\baselineskip}
\item{\cventry{2024}{STM32 Embedded State Machine}{Personal Project}{C, STM32 HAL, Embedded Systems}{}{\vspace{1pt}
\begin{itemize}
    \item Designed and implemented a timer-driven finite state machine in C on STM32 for multi-functional button input handling, demonstrating structured software design in resource-constrained environments
\end{itemize}}}

\end{itemize}

% ============================================================
%     PROFESSIONAL EXPERIENCE
% ============================================================

\section{Professional Experience}
\vspace{3pt}
\begin{itemize}

\needspace{5\baselineskip}
\item{\cventry{Apr 2026 - Jul 2026}{Research Assistant (10h/week)}{Munich University of Applied Sciences}{Munich, Germany}{}{\vspace{1pt}
\begin{itemize}
    \item Automated analysis and reporting of Sweep Frequency Response Analysis (SFRA) measurement data using a custom Python script, converting raw acquisition output into structured technical reports
    \item Designed and assembled electronic circuits and data acquisition setups for fault-detection experiments on rotating electrical machines
    \item Collaborated with researchers to iteratively improve testing procedures, communicating findings clearly across different levels of technical expertise
    \item Maintained detailed experimental documentation to ensure reproducibility and enable knowledge transfer within the team
\end{itemize}}}

\end{itemize}

% ============================================================
%     EDUCATION
% ============================================================

\section{Education}
\vspace{1pt}
\begin{itemize}

\needspace{5\baselineskip}
\item{\cventry{Oct 2024 - present}{Bachelor in Mechatronics}{Munich University of Applied Sciences (HM Munich)}{Munich, Germany}{}{\vspace{1pt}
In progress, expected 2027/2028. Core topics: embedded systems, control engineering, electronics, software engineering, mechanical design.
}}

\vspace{3pt}

\needspace{5\baselineskip}
\item{\cventry{Mar 2024 - Jul 2024}{Feststellungsprüfung, Technik (T/TI)}{Internationales Studienkolleg, Hochschule Kaiserslautern}{Kaiserslautern, Germany}{}{\vspace{1pt}
Overall grade: 1.5 --- Deutsch (2.0), Mathematik/Informatik (1.7), Physik (1.0). Equivalent to DSH-3/C1.
}}

\end{itemize}

% ============================================================
%     EXTRACURRICULAR
% ============================================================

\section{Extracurricular}
\vspace{3pt}
\begin{itemize}

\needspace{5\baselineskip}
\item{\cventry{Jan 2026 - present}{Active Member}{AkaModell München}{Munich, Germany}{}{\vspace{1pt}
\begin{itemize}
    \item Contributing to aircraft wing profile optimisation (MATLAB) and structural component design (Catia V5) in a student engineering team, applying analytical methods to real hardware development
\end{itemize}}}

\end{itemize}

% ============================================================
%     LANGUAGES
% ============================================================

\section{Languages}
\vspace{1pt}
\begin{itemize}
\item {Indonesian (native), English (fluent), German (fluent --- DSH-3/C1), Chinese (advanced).}
\end{itemize}

% ============================================================
%     REFERENCES
% ============================================================

\section{References}
\vspace{1pt}
\begin{itemize}
\item{Available upon request.}
\end{itemize}

\end{document}
